Метод Ньютона быстр только вблизи минимума: \hyperref[note:CDA139F9-5438-41E1-96B1-11A751BE40C5]{локально} он сходится квадратично, а \hyperref[note:CC41F7D2-EECB-46E4-817A-756E7CB9B7C0]{глобально} (с шагом Армихо--Вульфа, для сильно выпуклых функций) --- лишь как $O(1/\sqrt{k})$ по норме градиента. Вторая половина главы готовит ускорение без гессиана: \hyperref[note:866FE52A-109C-4912-B61C-F29BBD475F6D]{оценивающие последовательности} позволяют оценить невязку через $\lambda_k$ и строятся по явным формулам --- на них опирается \hyperref[note:BB5E3FC2-069C-44F0-92D0-B4DC13092CCA]{метод Нестерова}.
